%==============================================================================
% tento soubor pouzijte jako zaklad
% this file should be used as a base for the thesis
% (c) 2008 Michal Bidlo
% E-mail: bidlom AT fit vutbr cz
% Šablonu upravil / template edited by: Ing. Jaroslav Dytrych, dytrych@fit.vutbr.cz
%==============================================================================
% kodovaní: UTF-8 (zmena prikazem iconv, recode nebo cstocs)
% encoding: UTF-8 (you can change it by command iconv, recode or cstocs)
%------------------------------------------------------------------------------
% zpracování / processing: make, make pdf, make clean
%==============================================================================
% Soubory, které je nutné upravit: / Files which have to be edited:
%   xpolas34-ECS-20-literatura-bibliography.bib - literatura / bibliography
%   xpolas34-ECS-01-kapitoly-chapters.tex - obsah práce / the thesis content
%   xpolas34-ECS-30-prilohy-appendices.tex - přílohy / appendices
%==============================================================================
\documentclass[]{fitthesis} % bez zadání - pro začátek práce, aby nebyl problém s překladem
%\documentclass[english]{fitthesis} % without assignment - for the work start to avoid compilation problem
%\documentclass[zadani]{fitthesis} % odevzdani do wisu - odkazy jsou barevné
%\documentclass[english,zadani]{fitthesis} % for submission to the IS FIT - links are color
%\documentclass[zadani,print]{fitthesis} % pro tisk - odkazy jsou černé
%\documentclass[english,zadani,print]{fitthesis} % for the print - links are black
% * Je-li prace psana v anglickem jazyce, je zapotrebi u tridy pouzit 
%   parametr english nasledovne:
%   If thesis is written in english, it is necessary to use 
%   parameter english as follows:
%      \documentclass[english]{fitthesis}
% * Je-li prace psana ve slovenskem jazyce, je zapotrebi u tridy pouzit 
%   parametr slovak nasledovne:
%      \documentclass[slovak]{fitthesis}

% Základní balíčky jsou dole v souboru šablony fitthesis.cls
% Basic packages are at the bottom of template file fitthesis.cls
%zde muzeme vlozit vlastni balicky / you can place own packages here

%---rm---------------
\renewcommand{\rmdefault}{lmr}%zavede Latin Modern Roman jako rm / set Latin Modern Roman as rm
%---sf---------------
\renewcommand{\sfdefault}{qhv}%zavede TeX Gyre Heros jako sf
%---tt------------
\renewcommand{\ttdefault}{lmtt}% zavede Latin Modern tt jako tt

% vypne funkci šablony, která automaticky nahrazuje uvozovky,
% aby nebyly prováděny nevhodné náhrady v popisech API apod.
% disables function of the template which replaces quotation marks
% to avoid unnecessary replacements in the API descriptions etc.
\csdoublequotesoff

% =======================================================================
% balíček "hyperref" vytváří klikací odkazy v pdf, pokud tedy použijeme pdflatex
% problém je, že balíček hyperref musí být uveden jako poslední, takže nemůže
% být v šabloně
% "hyperref" package create clickable links in pdf if you are using pdflatex.
% Problem is that this package have to be introduced as the last one so it 
% can not be placed in the template file.
\ifWis
\ifx\pdfoutput\undefined % nejedeme pod pdflatexem / we are not using pdflatex
\else
  \usepackage{color}
  \usepackage[unicode,colorlinks,hyperindex,plainpages=false,pdftex]{hyperref}
  \definecolor{links}{rgb}{0.4,0.5,0}
  \definecolor{anchors}{rgb}{1,0,0}
  \def\AnchorColor{anchors}
  \def\LinkColor{links}
  \def\pdfBorderAttrs{/Border [0 0 0] }  % bez okrajů kolem odkazů / without margins around links
  \pdfcompresslevel=9
\fi
\else % pro tisk budou odkazy, na které se dá klikat, černé / for the print clickable links will be black
\ifx\pdfoutput\undefined % nejedeme pod pdflatexem / we are not using pdflatex
\else
  \usepackage{color}
  \usepackage[unicode,colorlinks,hyperindex,plainpages=false,pdftex,urlcolor=black,linkcolor=black,citecolor=black]{hyperref}
  \definecolor{links}{rgb}{0,0,0}
  \definecolor{anchors}{rgb}{0,0,0}
  \def\AnchorColor{anchors}
  \def\LinkColor{links}
  \def\pdfBorderAttrs{/Border [0 0 0] } % bez okrajů kolem odkazů / without margins around links
  \pdfcompresslevel=9
\fi
\fi
% Řešení problému, kdy klikací odkazy na obrázky vedou za obrázek
% This solves the problems with links which leads after the picture
\usepackage[all]{hypcap}

% Informace o práci/projektu / Information about the thesis
%---------------------------------------------------------------------------
\projectinfo{
  %Prace / Thesis
  project=BP,            %typ prace BP/SP/DP/DR  / thesis type (SP = term project)
  year=2017,             %rok odevzdání / year of submission
  date=\today,           %datum odevzdani / submission date
  %Nazev prace / thesis title
  title.cs={Komponentní systém pro herní grafický engine},  %nazev prace v cestine ci slovenstine (dle zadani) / thesis title in czech language (according to assignment)
  titleMain.cs={Komponentní systém\\ pro herní grafický engine},
  title.en={Graphic Engine Based on Entity Component System}, %nazev prace v anglictine / thesis title in english
  %Autor / Author
  titleMain.en={Graphic Engine Based on Entity Component System},
  author={Tomáš Polášek},   %cele jmeno a prijmeni autora / full name and surname of the author
  author.name={Tomáš},   %jmeno autora (pro citaci) / author name (for reference) 
  author.surname={Polášek},   %prijmeni autora (pro citaci) / author surname (for reference) 
  %author.title.p=Bc., %titul pred jmenem (nepovinne) / title before the name (optional)
  %author.title.a=PhD, %titul za jmenem (nepovinne) / title after the name (optional)
  %Ustav / Department
  department=UPGM, % doplnte prislusnou zkratku dle ustavu na zadani: UPSY/UIFS/UITS/UPGM
  %                  fill in appropriate abbreviation of the department according to assignment: UPSY/UIFS/UITS/UPGM
  %Skolitel / supervisor
  supervisor=Michal Španěl, %cele jmeno a prijmeni skolitele / full name and surname of the supervisor
  supervisor.name={Michal},   %jmeno skolitele (pro citaci) / supervisor name (for reference) 
  supervisor.surname={Španěl},   %prijmeni skolitele (pro citaci) / supervisor surname (for reference) 
  supervisor.title.p=,   %titul pred jmenem (nepovinne) / title before the name (optional)
  supervisor.title.a={Ing., Ph.D.},    %titul za jmenem (nepovinne) / title after the name (optional)
  %Klicova slova, abstrakty, prohlaseni a podekovani je mozne definovat 
  %bud pomoci nasledujicich parametru nebo pomoci vyhrazenych maker (viz dale)
  %Keywords, abstracts, declaration and acknowledgement can be defined by following 
  %parameters or using dedicated macros (see below)
  %===========================================================================
  %Klicova slova / keywords
  %keywords.cs={Klíčová slova v českém jazyce.}, %klicova slova v ceskem ci slovenskem jazyce
  %                                              keywords in czech or slovak language
  %keywords.en={Klíčová slova v anglickém jazyce.}, %klicova slova v anglickem jazyce / keywords in english
  %Abstract
  %abstract.cs={Výtah (abstrakt) práce v českém jazyce.}, % abstrakt v ceskem ci slovenskem jazyce
  %                                                         abstract in czech or slovak language
  %abstract.en={Výtah (abstrakt) práce v anglickém jazyce.}, % abstrakt v anglickem jazyce / abstract in english
  %Prohlaseni / Declaration
  %declaration={Prohlašuji, že jsem tuto bakalářskou práci vypracoval samostatně pod vedením pana ...},
  %Podekovani (nepovinne) / Acknowledgement (optional)
  %acknowledgment={Zde je možné uvést poděkování vedoucímu práce a těm, kteří poskytli odbornou pomoc.} % nepovinne
  %acknowledgment={Here it is possible to express thanks to the supervisor and to the people which provided professional help.} % optional
}

%Abstrakt (cesky, slovensky ci anglicky) / Abstract (in czech, slovak or english)
%\abstract[cs]{Cílem této práce je návrh a implementace knihovny pro správu entit tvořených pomocí kompozice z menších částí. Součástí práce je analýza dnes používaných technik pro práci s entitami -- \emph{objektově orientovaná dědičnost}, \emph{objektově orientovaná kompozice} a \emph{datově orientovaná kompozice} -- a jejich porovnání. \\
%Práce navrhuje komponentní systém pro jazyk C++, který je možno začlenit do existujících projektů. Vytvořené řešení umožňuje dynamickou kompozici entit za běhu aplikace z komponent definovaných za překladu. Systém je navržen s ohledem na snadnou paralelizaci a je uživatelsky rozšiřitelný. \\
%Výsledná knihovna umožňuje návrh aplikací za použití techniky \emph{datově orientované kompozice}.}
\abstract[cs]{Cílem této bakalářské práce je návrh a implementace knihovny pro správu entit, která umožňuje vývoj pomocí \emph{Entity-Component-System} paradigmatu. Součástí práce je analýza aktuálního stavu návrhu software, z pohledu vývoje her a rozbor dnes používaných technik pro práci s entitami včetně srovnání jejich výhod a nevýhod. Následně se práce zabývá důvody, proč tuto metodu návrhu používat a jakým způsobem využití kompozice ovlivňuje výkon aplikací, ve vztahu k hardwarové architektuře moderních počítačů, se zaměřením na paměťovou hierarchii.\\
V práci je dále navržen entitní systém, který umožňuje dynamickou kompozici entit za běhu aplikace z předem definovaných komponent. Výsledný systém je navržen s ohledem na snadnou paralelizaci vykonávaných akcí a umožňuje uživatelské rozšíření, ve směru datových struktur k uchovávání komponent. \\
Výsledkem této bakalářské práce je multiplatformní knihovna v programovacím jazyce \emph{C++}, pojmenovaná \emph{Entropy}, která vývojářům zpřístupňuje návrh aplikací pomocí \emph{ECS} paradigmatu.}
\abstract[en]{The goal of this bachelor thesis is the design and implementation of a library for entity management which allows software development using the \emph{Entity-Component-System} paradigm. Analysis of the current state of software design is presented, including comparison of techniques used for entity representation. Potential advantages of composition are discussed from the hardware point of view, primarily based on the memory hierarchy found in modern computers. \\
Thesis proposes design of an entity system, using the \emph{ECS} paradigm, which allows dynamic composition of entities during runtime of application from predefined components. The resulting system is designed to allow for easy parallelization of performed actions and user customizability of component data structures.\\
The result of this bachelor thesis is a \emph{C++} multi-platform library, named \emph{Entropy}, which enables its users to design applications using the \emph{ECS} paradigm.}

%Klicova slova (cesky, slovensky ci anglicky) / Keywords (in czech, slovak or english)
\keywords[cs]{Komponentní systém, Entitní systém, Entity-Component-System, ECS, Kompozice, Data-oriented design, Paralelismus, C++, Návrh her}
\keywords[en]{Component system, Entity-component system, Entity-component-system, ECS, Composition, Data-oriented design, Parallelism, C++, Game design}

%Prohlaseni (u anglicky psane prace anglicky, u slovensky psane prace slovensky)
%Declaration (for thesis in english should be in english)
\declaration{Prohlašuji, že jsem tuto bakalářskou práci vypracoval samostatně pod vedením pana Ing. Michala Španěla, Ph.D.
Uvedl jsem všechny literární prameny a publikace, ze kterých jsem čerpal.}

% \declaration{Hereby I declare that this bachelor's thesis was prepared as an original author’s work under the supervision of Mr. X
% The supplementary information was provided by Mr. Y
% All the relevant information sources, which were used during preparation of this thesis, are properly cited and included in the list of references.}

%Podekovani (nepovinne, nejlepe v jazyce prace) / Acknowledgement (optional, ideally in the language of the thesis)
\acknowledgment{Děkuji vedoucímu bakalářské práce, panu Ing. Michalu Španělovi, Ph.D. za odborné vedení a konzultace problémů, které se při její tvorbě vyskytly.}
%\acknowledgment{Here it is possible to express thanks to the supervisor and to the people which provided professional help
%(external submitter, consultant, etc.).}

% řeší první/poslední řádek odstavce na předchozí/následující stránce
% solves first/last row of the paragraph on the previous/next page
\clubpenalty=10000
\widowpenalty=10000

\usepackage[czech, linesnumbered, lined, algoruled, longend, norelsize]{algorithm2e}
\usepackage{algorithmicx}
\usepackage{algpseudocode}
\usepackage{subcaption}
\usepackage{cleveref}

\begin{document}
  % Vysazeni titulnich stran / Typesetting of the title pages
  % ----------------------------------------------
  \maketitle
  % Obsah
  % ----------------------------------------------
  \setcounter{tocdepth}{1}
  %\setcounter{secnumdepth}{1}
  \tableofcontents
  
  % Seznam obrazku a tabulek (pokud prace obsahuje velke mnozstvi obrazku, tak se to hodi)
  % List of figures and list of tables (if the thesis contains a lot of pictures, it is good)
\ifczech
  \renewcommand\listfigurename{Seznam obrázků}
\fi
\ifslovak
  \renewcommand\listfigurename{Zoznam obrázkov}
\fi

  % \listoffigures
\ifczech
  \renewcommand\listtablename{Seznam tabulek}
\fi
\ifslovak
  \renewcommand\listtablename{Zoznam tabuliek}
\fi

  % \listoftables 

  % vynechani stranky v oboustrannem rezimu
  % Skip the page in the two-sided mode
  \iftwoside
    \cleardoublepage
  \fi

  % Text prace / Thesis text
  % ----------------------------------------------
  %=========================================================================
% (c) Michal Bidlo, Bohuslav Křena, 2008

\iffalse

Způsoby, jak reprezentovat a uchovávat \emph{entity} existuje velké množství. V začátcích herního průmyslu byly \emph{entity} velmi úzce provázány s okolním kódem, ale díky velikosti týmů, které na hrách pracovaly, a relativní jednoduchosti her toto nebyl problém. Avšak s růstem složitosti se začaly používat další přístupy, jako například \emph{objektově orientované programování}, kdy je využíván \emph{polymorfismus} a \emph{dědičnost}. 

Cílem této práce je návrh a implementace knihovny pro dynamickou kompozici objektů, za běhu aplikace. Mezi hlavní požadavky patří \emph{thread-safe}\cite{ThreadSafety} rozhraní, které je možné používat ve vícevláknových aplikacích. Další prioritou je skutečné využití a prospěch z běhu na více vláknech, jednoduchost použití a integrace do již existujících projektů a minimalizace režie.

Práce je rozdělena do čtyř logických celků -- teoretická část\ref{Chap:Theory}, návrh systému\ref{Chap:Design}, implementace\ref{Chap:Implementation}, vyhodnocení a praktická použitelnost výsledné knihovny.

Kapitola \ref{Chap:Theory} se zabývá aktuálním stavem v návrhu \emph{software} a důvody proč je metoda \emph{datově orientované kompozice} důležitá. Jelikož efektivní využití architektury moderních počítačů je stále více důležité, byly při návrhu knihovny použity techniky datově orientovaného návrhu\cite{DataOrientedDesign}, jehož výhody a bližší popis je součástí kapitoly \ref{Chap:DDD}. Další důležitou částí je efektivní využití vícevláknových procesorů, kterým se zabývá kapitola \ref{Chap:Parallelism}.

Kapitola \ref{Chap:Design} obsahuje obecný návrh systému pro dynamickou kompozici, který není závislý na implementačním jazyce. První část \ref{Chap:Overview} obsahuje přehled celého systému, vymezení funkcionality jednotlivých bloků a jejich komunikace. Následují sekce o \emph{komponentech}(\ref{Chap:Component}), \emph{systémech}(\ref{Chap:System}) a \emph{entitách}(\ref{Chap:Entity}), ve kterých je specifikována jejich funkce. Návrhem paralelního přístupu a popis podporovaných možností paralelizace se zabývá sekce \ref{Chap:ParallelismDesign}.

Následuje kapitola \ref{Chap:Implementation}, ve které je popsána implementace výše zmíněného návrhu v jazyce \emph{C++}. První část se zabývá psaním přenositelného kódu v jazyce \emph{C++} a zdůvodnění výběru implementačního jazyka. Následuje popis implementace jednotlivých částí návrhu. Sekce \ref{Chap:CompImpl} obsahuje způsob registrace \emph{komponentů}, definuje, co může být \emph{komponentem} a základní typy datových struktur pro jejich uchovávání. Následuje popis \emph{systémů}, jejich definice a způsob, jakým jsou určeny k používání. Dále je zde shrnuta implementace \emph{entity} a výběr mezi horizontálními a vertikálními \emph{metadaty} s ohledem na paralelizmus. Nakonec je popsána funkce \emph{refresh} a způsoby zajištění parlelního přístupu.

Poslední část se skládá z příkladu využití knihovny v kapitole \ref{Chap:Demo} a vyhodnocení \ref{Chap:Eval}. Knihovna je zhodnocena z pohledu použitelnosti (\ref{Chap:Usability}), výkonosti (\ref{Chap:Performance}) a nakonec porovnána proti vybraným \emph{open-source} knihovnám.

Jelikož se využití \emph{komponentních systémů} začalo objevovat až v posledních letech, neexistuje obecně uznávaný způsob, jak by přesně měl vypadat. Existuje mnoho různých prací na toto téma

Právě tato metoda -- \emph{kompozice} místo \emph{dědičnosti} -- se začíná v herním průmyslu používat stále častěji. Díky celkové neprobádanosti tohoto návrhu existuje mnoho různých způsobů jak by takový systém měl vypadat a jaké části by měl obsahovat.

Tato metoda vzniká z použití tzv. \emph{datově orientovaného návrhu}\cite{DOD}, který se zabývá návrhem aplikací skrz analýzu a transformace dat. Jelikož je tento způsob návrhu stále vcelku nový, existuje mnoho způsobů, jak systém \emph{dynamické kompozice} navrhnout. 

Důležitou vlastností entitního systému založeného na kompozici je právě modularita, která umožňuje ...

Tato metoda vzniká z použití tzv. \emph{datově orientovaného návrhu}\cite{DOD}.

První dokumentované použití se objevuje ve hře \emph{Dungeon Siege}\cite{DungeonSiege} \emph{Tony Hawk}, kdy \emph{Mick West} během několika let transformoval celou hierarchii na právě kompozici z \emph{komponentů}. Následoval po

\fi

\chapter{Úvod}

Díky rostoucím požadavkům na moderní hry -- jak ve složitosti herních principů, tak i věrnosti grafické reprezentace -- se stále zvyšují požadavky na \emph{herní engine}, použitých v jejich tvorbě. Mezi hlavní požadavky patří široké využití v mnoha herních žánrech, možnosti multiplatformního nasazení, efektivní použití dostupného hardware, ale také jednoduchost a efektivita práce ve velkých týmech. Tyto, ale i další problémy řeší návrhový vzor \emph{entity-component-system}, který v základu používá \emph{kompozici} místo \emph{dědičnosti}.

Jelikož je tato metoda vcelku nová, širší využití začíná až v posledních letech, existuje mnoho návrhů, jak by \emph{entitní systém} založený na kompozici měl vypadat. Nejčastěji je návrh rozdělen do tří celků -- \emph{entita}, \emph{komponent} a \emph{systém}. Často se také využívá tzv. \emph{datově orientovaného návrhu}\cite{DOD}, který specifikuje návrh aplikace za použití analýzy vstupních a výstupních dat jednotlivých částí.

Cílem práce je návrh a implementace \emph{entitního systému} založeného na \emph{datově orientované kompozici}. Hlavními požadavky na tento systém jsou: 
\begin{itemize}
	\item Použitelnost ve vícevláknových aplikacích.
	\item Efektivní využití hardwarových prostředků.
	\item Jednoduché rozhraní a integrace do existujících projektů.
\end{itemize}
Mezi další požadavky, z pohledu \emph{herního enginu}, jsou efektivní práce s velkým množstvím \emph{entit}, kde většina z nich nemusí být aktuálně používána a také komunikace mezi jednotlivými \emph{komponenty} a \emph{entitami}.

Práce je rozdělena do čtyř logických celků -- teoretická část, návrh systému, implementace a vyhodnocení. Kapitola \ref{Chap:Theory} popisuje aktuální stav návrhu \emph{software}, primárně z pohledu herního vývoje. Nejdříve zde jsou přiblíženy požadavky na \emph{entitní systémy}. Následuje porovnání aktuálně používaných způsobů reprezentace \emph{entit}, jejich výhody a nevýhody. Závěr obsahuje způsoby využití vícevláknových konfigurací. Kapitola \ref{Chap:Design} zahrnuje kompletní návrh \emph{entitního systému}, který je obecný -- nezávislý na implementačním jazyce. Nejdříve je uveden přehled celého systému, následuje podrobnější popis jeho částí a rozbor různých způsobů řešení. Další kapitola (\ref{Chap:Implementation}) je věnována implementaci v jazyce C++. Následuje kapitola \ref{Chap:Demo}, ve které je předvedena funkčnost systému na návrhu a implementaci demo hry. Poslední kapitola shrnuje výsledný systém z pohledu výkonosti a srovnává jej s podobnými \emph{open-source} knihovnami.

\chapter{Teoretický rozbor}

Tato se věnuje teorii \emph{entitních systémů}, jejich historii a požadavkům na ně kladeným. Dále obsahuje rozbor běžně používaných způsobů reprezentace \emph{entit} a srovnává jejich výhody a nevýhody. Následuje popis vlastního \emph{entitního systému} založeného na kompozici, definice hlavních pojmů -- \emph{entita}, \emph{komponent} a \emph{systém}. Na závěr se věnuje metodám paralelizmu, které jsou dnes používány.

\section{Entitní systémy a jejich historie}
\todo{Zaklad, co to je.}
\blind[1]

Entitní systém je část, v tomto případě, \emph{herního enginu}, jejíž úkolem je údržba \emph{entit}, kde \emph{entitou} je objekt ve virtuálním světě. 

\blind[3]

\section{Požadavky na návrh ####}
\todo{Délka, složitost. Potřeba oddělení částí.}
\blind[2]

Stále zvyšující se požadavky na složitost herních principů, věrnost grafické reprezentace a velikost virtuálních světů mnohonásobně ztěžují návrh \emph{herních enginů}, na kterých jsou hry stavěny. Dalším problémem je znovupoužitelnost již vytvořených částí, nejen ve stejném projektu, ale stále častěji i v dalších hrách. Z výše zmíněných důvodů vznikají techniky pro organizaci kódu a knihy \emph{návrhových vzorů}\cite{DesignPatterns}\cite{GameDesignPatterns}, které obsahují zkušenosti a prověřené způsoby jak navrhovat \emph{software}. Důležitou vlastností správného návrhu, je také modularita -- oddělení částí systému. Tyto techniky umožňují pracovat velkému množství programátorů na jednom projektu.

Mezi techniky dnes používané patří primárně \emph{objektově orientované programování}, ale stále častěji také \emph{}

Způsob reprezentace \emph{entit} -- objektů ve virtuálním světě -- je jedním z důležitých rozhodnutí v návrhu \emph{herního enginu}. Entity jsou často používány pro komunikaci mezi jednotlivými podsystémy -- např. vykreslení efektu, který vznikl v \emph{gameplay} kódu. Přílišná provázanost však většinou znamená zvýšený výskyt programovacích chyb (\emph{bugs}) \cite{GameDesignPatterns}. Podrobněji se tímto zabývá sekce \ref{Chap:Representation}.

\section{Datově orientovaný návrh}
\label{Chap:DDD}

Dnes nejpoužívanější způsob návrhu -- \emph{objektově orientovaný} (\emph{OOD}) -- umožňuje abstrahovat od fyzického hardware, na kterém výsledná aplikace běží a řešit daný problém jeho dekompozicí do objektů. Základní stavební kámen je v tomto případě objekt -- agregace hodnot a proveditelných operací. Toto umožňuje řešení problémů transformací objektů z reálného světa na objekty virtuální, se kterými dokáží lidé pracovat a zároveň jim rozumí také počítač. Toto je však problematické pro moderní výpočetní hardware, který je optimalizovaný na jednoduché datové transformace.

Kvůli výše zmíněnému problému a stále vyšší potřebě optimalizace se objevuje \emph{datově orientovaný návrh}\cite{DOD} (\emph{DOD}), který se naopak zaměřuje na data, které aplikace používá. Nevýhodou je snížení čitelnosti výsledného kódu a složitější transformace návrhu aplikace -- řešení problému -- ve výsledný program. Základní myšlenkou je oddělení dat a operací nad nimi, toto umožňuje efektivnější využití paměti a procesorových \emph{cache}. 

Častým příkladem rozdílů mezi \emph{OOD} a \emph{DOD} je transformace seznamu struktur (\emph{SOA}) ve strukturu seznamů (\emph{AOS}), který lze vidět na obr.\ref{Fig:SOAASO}.

\begin{figure}[]
	\tmpframe{\includegraphics[width=\linewidth]{TODO-image}}
	\caption{Příklad transformace pole struktur na strukturu polí. \todo{Vytvořit}}
	\label{Fig:SOAASO}
\end{figure}

Základem \emph{DOD} je důkladná analýza aplikačních dat -- jejich obor hodnot, transformace, definice vstupních dat, požadovaný výstup atp. Často je využíváno datové struktury pole, jehož vlastnosti umožňují rychlou iteraci. Mezi výhody \emph{DOD} patří také dobrá lokalita dat a odkazů, což umožňuje efektivní využití procesorových \emph{cache}. Díky těmto vlastnostem se \emph{DOD} stále častěji používá v herním průmyslu\cite{DataOrientedDesignDice}\cite{DataOrientedDesignCppCon}.

Entitní systém navržený v této práci vychází z principů \emph{DOD} -- \emph{datově orientovaná kompozice}\cite{DODComponents}. Blíže se tímto způsobem skládání objektů zabývá sekce \ref{Chap:DOC}.

\section{Reprezentace entit}

Tato sekce obsahuje rozbor nejpoužívanějších způsobů reprezentace \emph{entit}\cite{EvolveHierarchy} a jejich \emph{chování}. Pod pojmem \emph{entita} je myšlen objekt ve virtuálním světě a \emph{chování} jsou operace, které \emph{entita} dokáže provést -- pohyb, vykreslení, kolize atp.

Jednotlivé reprezentace jsou srovnány na příkladu jednoduché 2D \emph{shoot'em up} hry. Metody jsou hodnoceny na základě složitosti implementace, návrhu s jejich použitím a výhodnosti z pohledu hardware. 

\subsection{Objektově orientovaná hierarchie}
\todo{Popis, příklad, vyhodnocení výhod/nevýhod, komunikace}
\begin{figure}[H]
	\tmpframe{\includegraphics[width=\linewidth]{TODO-image}}
	\caption{Příklad objektově orientované hierarchie. \todo{Vytvořit}}
\end{figure}
\begin{figure}[H]
	\tmpframe{\includegraphics[width=\linewidth]{TODO-image}}
	\caption{Obsah paměti u objektově orientovaných hierarchií. \todo{Vytvořit}}
\end{figure}
\blind[3]

 problem nastava v pripade 
 nekolikanasobna dedicnost, problemy

\subsection{Objektově orientovaná kompozice}
\todo{Popis, příklad, vyhodnocení výhod/nevýhod, komunikace}
\begin{figure}[H]
	\tmpframe{\includegraphics[width=\linewidth]{TODO-image}}
	\caption{Příklad objektově orientované kompozice. \todo{Vytvořit}}
\end{figure}
\begin{figure}[H]
	\tmpframe{\includegraphics[width=\linewidth]{TODO-image}}
	\caption{Obsah paměti u objektově orientované kompozice. \todo{Vytvořit}}
\end{figure}
\blind[3]

\subsection{Datově orientovaná kompozice}
\label{Chap:DOC}
\todo{Popis, příklad, vyhodnocení výhod/nevýhod, komunikace}
\begin{figure}[H]
	\tmpframe{\includegraphics[width=\linewidth]{TODO-image}}
	\caption{Příklad datově orientované kompozice. \todo{Vytvořit}}
\end{figure}
\begin{figure}[H]
	\tmpframe{\includegraphics[width=\linewidth]{TODO-image}}
	\caption{Obsah paměti u datově orientované kompozice. \todo{Vytvořit}}
\end{figure}
\blind[4]

\section{Entitní systém}
\subsection{Koncept}
\todo{Představení základní architektury návrhového vzoru - entity, komponenty, systémy}
\begin{figure}[H]
	\tmpframe{\includegraphics[width=\linewidth]{TODO-image}}
	\caption{Entita reprezentována jako řádek v databázi. \todo{Vytvořit}}
\end{figure}
\blind[5]

\subsection{Výhody zvolené reprezentace}
\todo{Lokalita dat, instrukcí. Cache. Dynamičnost, prototypování. Serializace. Embedded skriptovací jazyky.}
\blind[4]

\section{Paralelizmus}
\label{Chap:Parallelism}
\todo{Důležitost, hardware - cache, popis možností paralelizace.}
\blind[2]
\begin{figure}[H]
	\tmpframe{\includegraphics[width=\linewidth]{TODO-image}}
	\caption{Paralelizmus uvnitř systémů. \todo{Vytvořit}}
\end{figure}
\blind[1]
\begin{figure}[H]
	\tmpframe{\includegraphics[width=\linewidth]{TODO-image}}
	\caption{Paralelizmus na úrovni systémů. \todo{Vytvořit}}
\end{figure}
\blind[1]
%https://software.intel.com/en-us/articles/designing-the-framework-of-a-parallel-game-engine

\chapter{Návrh entitního systému}
\todo{Co by to melo umet, priority}

\secion{Přehled a komunikace}
\todo{Postup seznámení s návrhem top -> down}
\begin{figure}[H]
	\tmpframe{\includegraphics[width=\linewidth]{TODO-image}}
	\caption{Návrhový diagram komponentního systému. \todo{Vytvořit}}
\end{figure}
\blind[1]
\todo{High-level design, propojení. Universe, action-cache.}
\begin{figure}[H]
	\tmpframe{\includegraphics[width=\linewidth]{TODO-image}}
	\caption{Přehled návrhu komponentního systému a komunikace mezi jeho částmi. \todo{Vytvořit - HLB diagram}}
\end{figure}
\blind[3]

\section{Komponenty a jejich nosiče}
\todo{Komponenta POD. Trivial constructor, copy. Paměť a typy nosičů. Typy komponent - tag, normal.}
\blind[4]

\section{Systémy a skupiny}
\todo{Volnost použití systémů. Uživatelem definované chování. Filtry - require, reject. Cache pomocí skupin entit.}
\blind[5]

\section{Reprezentace entit}
\todo{Entita, id, generace. Alternativy a důvody.}
\blind[3]

\section{Paralelní přístup}
\todo{Podporované způsoby paralelizace, návrh, zdůvodnění. Entity paralellization. System parallelization. Diskuze na téma, zda je paralelní přístup potřeba.}
\blind[5]

\section{Tok řízení}
\todo{Popsat předpokládaný tok řízení, jak se bude používat. Registrace. Iterace. Refresh.}
\begin{figure}[H]
	\tmpframe{\includegraphics[width=\linewidth]{TODO-image}}
	\caption{Digram toku řízení komponentního systému. \todo{Vytvořit - Control flow diagram}}
\end{figure}
\blind[3]


\chapter{Implementace}
\todo{Shrnutí návrhu}
\blind[2]

\section{Použité knihovny a přenositelnost}
\todo{zduvodneni pouziti c++ Zdůvodnění nepoužití knihoven. C++14, Linux, Windows. Header-only}
\blind[2]

\section{Komponenty}
\todo{Rekapitulace vlastností, které komponenta musi splnit. Registrace, generování identifikátorů. Uživatelské nosiče}
\blind[3]

\section{Systémy a skupiny}
\todo{Super třída, filtry, přístup - iterace. Přidání systému, volitelná skupina - added/removed.}
\blind[3]

\section{Správa entit}
\todo{Potřebné informace, implementace, datové struktury, optimalizace}
\blind[3]

\section{Obnovovací fáze}
\todo{Pořadí refresh. Obnovení aktuality skupin.}
\blind[2]

\section{Podpora paralelizmu}
\todo{2 způsoby paralelizmu - dělení skupin, change-sety}
\blind[3]


\chapter{Použití knihovny}

\section{Demo hra}
\todo{Popis principů hry}
\blind[1]

\section{Návrh}
\todo{Komponentní návrh}
\blind[3]

\section{Zhodnocení}
\todo{Výkon, návrh, jednoduchost}
\blind[3]

\chapter{Vyhodnocení}
\todo{Shrnutí, popis co bude v této kapitole}
\blind[1]

\section{Testování knihovny}
\todo{Unit testing, benchmarky}
\blind[2]

\section{Testované sestavy}
\todo{Hardware sestavy - PC, notebook. Tabulka identifikace}
\blind[1]

\section{Výkonnostní testy}
\todo{Výkon při různých zatíženích - procenta přidaných/odebraných. Grafy.}
\blind[1]
\todo{Porovnání nosičů komponent. Grafy.}
\blind[1]
\todo{Seriové vs paralelní - počet vláken. Grafy.}
\blind[1]

\section{Porovnání}
\todo{Výkonostní porovnání oproti jiným ECS knihovnám, Unity? Grafy.}
\blind[2]

\subsection{Implementace}
\todo{Použité knihovny, základní popis implementace}
\blind[2]

\subsection{Výsledek}
\todo{Poznatky a výsledky získané ze hry. Grafy?}
\blind[2]


\chapter{Závěr}
\todo{Využití komponentních systémů a aktuální stav} 
\blind[1]
\todo{Návrh, prototypy, priority} 
\blind[1]
\todo{Implementace, rozšiřitelnost, přenositelnost}
\blind[1]
\todo{Experimenty, srovnání, demo hra}
\blind[1]
\todo{Možnosti rozšíření}
\blind[1]

%\begin{figure}
%	\tmpframe{\includegraphics[width=\linewidth]{TODO-image}}
%	\caption{Obrazek \todo{Obrazek}}
%\end{figure}

%=========================================================================
 % viz. obsah.tex / see obsah.tex

  % Pouzita literatura / Bibliography
  % ----------------------------------------------
\ifslovak
  \makeatletter
  \def\@openbib@code{\addcontentsline{toc}{chapter}{Literatúra}}
  \makeatother
  \bibliographystyle{bib-styles/czechiso}
\else
  \ifczech
    \makeatletter
    \def\@openbib@code{\addcontentsline{toc}{chapter}{Literatura}}
    \makeatother
    \bibliographystyle{bib-styles/czechiso}
  \else 
    \makeatletter
    \def\@openbib@code{\addcontentsline{toc}{chapter}{Bibliography}}
    \makeatother
    \bibliographystyle{bib-styles/englishiso}
  %  \bibliographystyle{alpha}
  \fi
\fi
  \begin{flushleft}
  \bibliography{xpolas34-ECS-20-literatura-bibliography}
  \end{flushleft}

  % vynechani stranky v oboustrannem rezimu
  % Skip the page in the two-sided mode
  \iftwoside
    \cleardoublepage
  \fi

  % Prilohy / Appendices
  % ---------------------------------------------
  \appendix
\ifczech
  \renewcommand{\appendixpagename}{Přílohy}
  \renewcommand{\appendixtocname}{Přílohy}
  \renewcommand{\appendixname}{Příloha}
\fi
\ifslovak
  \renewcommand{\appendixpagename}{Prílohy}
  \renewcommand{\appendixtocname}{Prílohy}
  \renewcommand{\appendixname}{Príloha}
\fi
  \appendixpage

% vynechani stranky v oboustrannem rezimu
% Skip the page in the two-sided mode
\iftwoside
  \cleardoublepage
\fi
  
\ifslovak
%  \section*{Zoznam príloh}
%  \addcontentsline{toc}{section}{Zoznam príloh}
\else
  \ifczech
%    \section*{Seznam příloh}
%    \addcontentsline{toc}{section}{Seznam příloh}
  \else
%    \section*{List of Appendices}
%    \addcontentsline{toc}{section}{List of Appendices}
  \fi
\fi
  \startcontents[chapters]
  % seznam příloh / list of appendices
  % \printcontents[chapters]{l}{0}{\setcounter{tocdepth}{2}}
  
  % vynechani stranky v oboustrannem rezimu
  \iftwoside
    \cleardoublepage
  \fi
  % Tento soubor nahraďte vlastním souborem s přílohami (nadpisy níže jsou pouze pro příklad)
% This file should be replaced with your file with an appendices (headings below are examples only)

% Umístění obsahu paměťového média do příloh je vhodné konzultovat s vedoucím
% Placing of table of contents of the memory media here should be consulted with a supervisor
\chapter{Obsah přiloženého paměťového média}

\chapter{Manuál}

%\chapter{Konfigurační soubor} % Configuration file

%\chapter{RelaxNG Schéma konfiguračního souboru} % Scheme of RelaxNG configuration file

%\chapter{Plakát} % poster
 % viz. prilohy.tex / see prilohy.tex
\end{document}
